\section{Garantías de despliegue y reproducibilidad}

\subsection{Checklist de preparación para el despliegue}

Antes del despliegue es obligatorio revisar la prompt stability, el token budget, el plan de fallback y los hooks de monitoreo. Se recomienda usar una multi-column table para mostrar el estado actual de cada indicador y escribir el artifact de cada revisión en `deploy/logs/`.

\begin{table}[H]
\centering
\begin{tabular}{p{4cm}p{5cm}p{3cm}}
\toprule
Elemento de revisión & Método de verificación & Estado \\
\midrule
Prompt stability & múltiples prompts + CoT + self-consistency & En monitoreo \\
Token budget & latency + FLOPs shell & Dimensiones alineadas \\
Fallback & offline fallback script + human review & Preconfigurado \\
Monitoring hook & loss/latency/attention drift & hook enabled \\
\bottomrule
\end{tabular}
\caption{Checklist previa al despliegue.}
\end{table}

\begin{importantbox}{Uso de la checklist de preparación}
Antes de cada despliegue se recorre una standardized template: prompt log → evaluation dashboard → fallback script. Si algún elemento no se aprueba, debe registrarse en el incident grade-book y se pospone la puesta en producción.
\end{importantbox}

\subsection{Tooling stack de observabilidad y reproducibilidad}

Al construir el tooling stack se deben incluir la instrumentation, el seguimiento del LLM plan, la trace de las llamadas a herramientas (doc search, sondeo del python executor) y el attention/scoreboard snapshot.

\begin{figure}[H]
\centering
\includegraphics[width=0.85\textwidth]{slides-images/slide-37.jpg}
\caption{Deployment + stack de observabilidad (slide-37).}
\end{figure}

\begin{table}[H]
\centering
\begin{tabular}{p{3.5cm}p{5cm}}
\toprule
Nivel & Herramienta/práctica \\
\midrule
Prompt & Prompt log + versioned dashboard \\
Inference & Monitoring hook + latency/failure metrics \\
Tooling & API call trace + sandbox constraints \\
Observabilidad & Attention/Config snapshot + drift alert \\
\bottomrule
\end{tabular}
\caption{Stack de observabilidad del despliegue.}
\end{table}

\begin{warningbox}{Riesgos de carecer de observabilidad}
\begin{itemize}
    \item Sin attention trace, no es posible confirmar si un emergent jump es real;
    \item sin tool call log, las low-level skills no pueden rastrearse;
    \item la falta de drift alert retrasará el incident trigger.
\end{itemize}
\end{warningbox}

\subsection{Resumen de la sección}

La parte de despliegue y reproducibilidad es la barrera de protección de la práctica: desde la checklist hasta el stack de observabilidad y, después, el incident grade-book, el aim es que los high-level insights puedan llevarse a la práctica.

